\scanpage{85}\mspage{29}
% L1-2
Considérons
\[
 p\Phi_0:\mathrm{Cat}_{/S}\unsure{\mathord{-}|\longrightarrow}\mathrm{Cat}_{/S}
\]
% L3-5
ça transforme donc $X$ en $\Phi C(X)$, regardé comme objet de $\mathrm{Cat}_{/\deleted{\unsure{s}}S}$ sans plus. On a un morphisme fonctoriel
% L6
\[
 X\xrightarrow{j_X} C(X)\qquad j:\mathrm{id}_{\mathrm{Cat}_{/S}}\longrightarrow C
\]
% L7-8
d'où un morphisme fonctoriel
\[
 \Phi(X)\xrightarrow{\Phi(j_X)}\Phi C(\deleted{\unsure{x}}X)\qquad
 \Phi*j:\Phi\longrightarrow\Phi\circ C,
\]
% L9-10
et également un morphisme fonctoriel
\[
 X\xrightarrow{i_X}\Phi(X)\deleted{\unsure{\longrightarrow\Phi C(X)}}
\]
% L11-12
d'où en composant
\diagram{\begin{tikzcd}[ampersand replacement=\&]
 X\arrow[r,"i_X"]\arrow[rr,bend right=28,"\alpha_X"']\&\Phi(X)\arrow[r,"{\Phi(j_X)}"]\&\Phi C(X),
\end{tikzcd}}
% L13-14
\unsure{i.e.|ou}
\[
 \alpha:\mathrm{id}_{\mathrm{Cat}_{/S}}\longrightarrow p\Phi_0
\]
% L15-19
Pour $X$ fixé dans $\mathrm{Cat}_{/S}$, on sait que $j_X$ est dans $W$, donc aussi $\Phi(j_{\unsure{x|X}})$ (quoique non fibre par fibre), cf. cor.~1 \unsure{page|p.}~20. On sait aussi que $i_X$ est dans $W$, donc aussi \deleted{\unsure{le}} le composé $\alpha_X$.

% L20-21
Considérons
\[
 \Phi_0 p:\underline{\mathrm{Fibsc}}_0\,S\longrightarrow\underline{\mathrm{Fibsc}}_0\,S.
\]
% L22-23
On part de $X$ dans $\underline{\mathrm{Fibsc}}_0\,S$, on oublie le scindage et le fait même qu'il soit \unsure{Hot-fibrant|Cat-fibrant}
% archive sheet number: 84
